% Generated draft: requires mathematical and source review.
\documentclass[11pt]{article}
\usepackage[margin=25mm]{geometry}
\usepackage{fontspec}
\setmainfont{Latin Modern Roman}
\usepackage{amsmath,amssymb,mathtools,unicode-math}
\setmathfont{Latin Modern Math}
\usepackage{enumitem,needspace,xurl,hyperref}
\setlength{\parindent}{0pt}
\setlength{\parskip}{5pt}
\setlength{\emergencystretch}{3em}
\newcommand{\N}{\mathbb{N}}
\newcommand{\Z}{\mathbb{Z}}
\newcommand{\Q}{\mathbb{Q}}
\newcommand{\R}{\mathbb{R}}
\newcommand{\C}{\mathbb{C}}
\newcommand{\F}{\mathbb{F}}
\newcommand{\A}{\mathbb{A}}
\newcommand{\PP}{\mathbb{P}}
\newcommand{\E}{\mathbb{E}}
\newcommand{\eps}{\varepsilon}
\newcommand{\dd}{\,\mathrm{d}}
\begin{document}
\section*{UnsolvedMath ID: 30000378}
\textbf{Source problem code:} catalogue entry
\subsection{Definitions and mathematical statement}
Let $H$ be a Borel group in the sense intended by the report, and write $\tau_0$ for its ``original topology.'' The extraction does not specify whether $H$ is a Borel subgroup of an ambient Polish group with its subspace topology, a standard Borel group with additional topological data, or another convention. This distinction must be resolved before the topological condition is definite.

Let $\mathcal T(H,\tau_0)$ be the class of Hausdorff locally compact second-countable group topologies $\tau$ on the same underlying group such that $\tau_0\subseteq\tau$. The inclusion expresses the source's word ``extending'' as refinement. Whether the Borel structure must also remain unchanged is not stated and is not silently imposed.

Define $E_0$ on $2^{\N}$ by eventual equality. On $(2^{\N})^{\N}$ define its countable product by
\[
 x\ E_0^{\N}\ y
 \quad\Longleftrightarrow\quad
 \forall j\in\N\ \exists N_j\in\N\ \forall n\ge N_j:
       x(j)(n)=y(j)(n).
\]
Each coordinate may have its own eventual-equality threshold; no common threshold across all coordinates is required.

For a Polish $H$-space $X$ with its intended action, the orbit relation is
\[
 u\ E_H^X\ v\quad\Longleftrightarrow\quad
 \exists h\in H:\ h\cdot u=v.
\]
A Borel reduction of $E_0^{\N}$ to $E_H^X$ is a Borel map $f:(2^{\N})^{\N}\to X$ satisfying
\[
 x\ E_0^{\N}\ y\quad\Longleftrightarrow\quad
 f(x)\ E_H^X\ f(y)
 \qquad\text{for all }x,y.
\]
The question asks whether, for every $H$ in the intended Borel-group class,
\[
 \mathcal T(H,\tau_0)=\varnothing
 \quad\Longleftrightarrow\quad
 \exists\text{ Polish }H\text{-space }X:
             E_0^{\N}\le_B E_H^X.
\]
The usual term Polish $H$-space means a continuous action on a separable completely metrizable space, but the topology or Borel-action convention for $H$ is not specified here. No requirement of freeness, transitivity, or a particular action is added. The question concerns the absence of an extending locally compact topology, not merely the failure of $\tau_0$ itself to be locally compact.
\subsection{Short English statement}
The proposed equivalence links the impossibility of refining the original group topology to a locally compact second-countable one with the existence of an action whose orbit relation contains coordinatewise eventual equality in Borel complexity. The meaning of Borel group and the action topology need clarification.
\subsection{Sources}
\begin{itemize}
\item Local UnsolvedMath registry, permanent ID 30000378, source code catalogue entry, ``Orbit Equivalence Relations of Nonlocally Compact Borel Groups.'' Proposed year 2005; proposer unspecified.
\item \emph{Set Theory} (2006), Oberwolfach Reports, pp.\ 3121--3174, DOI 10.4171/owr/2005/55; \url{https://doi.org/10.4171/owr/2005/55}. The displayed report year 2006 differs from the DOI and proposed year 2005.
\item Supplied reference labeled ``later partial locally-compact converse'': \url{https://www.impan.pl/shop/publication/transaction/download/product/91329}. No full bibliographic identity or theorem hypotheses are supplied.
\item Statement status corrected\_verified; the supplied verification says group, relation, and sequence-space notation were restored without changing the equivalence or definition.
\item Literature triage checked 2026-08-21: open. It says a compact case is solved by an $E_0$ criterion while the displayed $E_0^{\N}$ locally compact criterion is open for general Borel groups, with important subclasses known. Exact scope and current status require independent verification. Created 2026-08-25; updated 2026-09-16.
\item Registry curation by Ulam AI / UnsolvedMath Contributors, CC BY 4.0; underlying report and literature provenance remain as supplied. Unreviewed statement-only draft.
\end{itemize}
\end{document}
